% chapters/journal/02-related-work.tex
% 第2章 相关工作

\section{相关工作}

\subsection{传统交易排序策略}

以太坊等区块链平台的默认交易排序遵循 Gas Price 贪心规则：区块构建者按交易手续费从高到低排列候选交易，在区块 Gas 上限约束下依次打包\cite{ethereum2024pbs}。该策略的优化目标单一，仅关注短期手续费收益，对于交易间是否存在状态冲突、执行顺序是否合理等问题不做判断。以太坊 EIP-1559 引入基础费用和优先费用的双层定价机制后，交易手续费的结构更加复杂，简单按优先费用排序并不总能反映交易的真实价值和紧迫程度。

先进先出（FIFO）排序按交易到达时间顺序打包，在公平性方面优于贪心排序，但完全放弃了对收益的优化。\"Oz 等\cite{oz2024fcfs}在先到先得排序机制下研究了 MEV 提取行为，指出即使采用时序排列，攻击者仍可通过网络层面的时延操纵影响交易顺序，FIFO 并不能从根本上消除排序操纵的可能性。Hay 等\cite{hay2025fairness}提出基于乐观策略的交易排序公平性协议，试图在去中心化环境下保证交易排序的时序公平，但其依赖可信时间源假设，在实际部署中面临时钟同步和网络延迟等工程挑战。上述方法均为基于静态规则的排序策略，在设计时预设了固定的排序准则，缺乏对交易间依赖关系和排序风险的动态感知能力。当候选交易集合的构成发生变化时，静态规则无法做出自适应调整。

\subsection{MEV 问题与应对方案}

Daian 等\cite{daian2019flash}在 Flash Boys 2.0 中系统揭示了以太坊 DEX 场景下的 MEV 问题，指出套利机器人通过抢跑和三明治攻击操纵交易排序，从普通用户交易中获取额外收益。Mazorra 等\cite{mazorra2022price}从博弈论视角分析了 MEV 的价格效应。在工程实践方面，Flashbots 提出 MEV-Boost\cite{flashbots2022mevboost}，通过中继将验证者与构建者解耦，但该机制不改变构建者内部的排序逻辑。Vedula 等\cite{vedula2023masquerade}提出 Masquerade 方案，通过轻量级交易重排缓解 MEV，但其重排规则为静态设计，适应性有限。Alipanahloo 等\cite{alipanahloo2024mev}对现有 MEV 缓解方案进行了全面综述，指出在构建者层面缺乏同时兼顾收益与风险控制的排序方法。

\subsection{强化学习在交易排序中的应用}

强化学习在组合优化和序列决策问题上已展现出良好性能\cite{sutton2018rl}。区块构建本质上是在资源约束下从候选集合中构造有序序列的问题：给定一组候选交易和区块 Gas 容量上限，目标是选择交易子集并确定执行顺序，使得特定目标函数最大化。该问题与经典的带约束排列优化具有相似的数学结构，而强化学习的序列决策框架天然适合对这类问题进行建模。

在区块链交易排序领域，已有若干将深度强化学习引入的尝试。Wen 等\cite{wen2025rlbes}提出 RL-BES，采用双 DRL 网络结合蒙特卡洛树搜索进行区块构建，在提升 MEV 提取效率方面取得了效果，但其优化目标侧重于构建者收益最大化，未考虑用户公平性和风险抑制。Seoev 等\cite{seoev2025bidding}将强化学习用于 Polygon 链上的 MEV 提取竞价策略，关注点在于竞价博弈而非排序本身。在区块链其他领域，深度强化学习也被用于共识参数优化\cite{zou2023consensus}和支付通道调度\cite{ren2024payment}等任务，证明了 RL 在区块链系统优化中的普适性。

表~\ref{tab:comparison} 从动作空间设计、优化目标和风险考量三个维度对比了本文方法与已有 RL 排序工作的区别。

\begin{table}[htbp]
\centering
\caption{本文方法与已有 RL 排序方法的对比}
\label{tab:comparison}
\begin{tabular}{lccc}
\toprule
方法 & 动作空间 & 优化目标 & 风险考量 \\
\midrule
RL-BES\cite{wen2025rlbes}  & 规则切换 & MEV 提取效率 & 无 \\
DQN 排序     & 规则切换 & 综合指标 & 有限 \\
本文方法     & 逐笔选交易 & 收益+公平+风险 & 复合奖励 \\
\bottomrule
\end{tabular}
\end{table}

与上述工作相比，本文方法在动作空间设计上具有本质区别。现有 RL 排序方法的动作空间通常定义为在 FIFO、Gas 优先等少量预设规则之间切换，智能体每步选择一种排序规则应用于整个交易集合。这种"选规则"的动作设计将排序决策简化为分类问题，策略网络只能在有限的规则集合中做出选择，无法刻画交易间的位置依赖关系。本文将动作空间定义为"逐笔选交易"：智能体在每步从当前候选集合中选择一笔具体交易，经过多步决策生成完整的排序序列。该设计使策略网络能够感知每笔交易在序列中的位置效应，根据已选交易的构成动态调整后续选择，从而在更细的粒度上优化排序结果。
